% ===========================================================================
% Overleaf-ready technical note — grid-forming converter model choice for VIS.
% Self-contained (manual bibliography); compile: pdflatex model_choice (x2).
% ===========================================================================
\documentclass[11pt,a4paper]{article}
\usepackage[margin=1in]{geometry}
\usepackage{amsmath,amssymb}
\usepackage{booktabs}
\usepackage{siunitx}
\usepackage{xcolor}
\usepackage[colorlinks=true,allcolors=blue]{hyperref}
\usepackage{cleveref}
\usepackage{enumitem}

\title{\textbf{Which ANDES Converter Model for Virtual Inertia Scheduling?}\\[2pt]
\large REGCV1 vs.\ the REGF Grid-Forming Family --- a Model-Selection Note}
\author{Mehdi Abdallahi \\ \small RRE, ETH Zurich / EPFL}
\date{\today}

\newcommand{\model}[1]{\texttt{#1}}

\begin{document}
\maketitle

\begin{abstract}
\noindent Virtual inertia scheduling (VIS) co-optimises generation dispatch with
the \emph{virtual inertia} $M$ and \emph{damping} $D$ provided by grid-forming
inverters. The choice of converter dynamic model therefore determines whether the
scheduler's decision variables even exist in the simulator. This note compares the
seven renewable/converter models in ANDES~1.9.3 and concludes---from the model
equations, the physical (swing-equation) requirement, and comparability with the
foundational VIS literature---that \textbf{\model{REGCV1}} is the correct model for
this study. \model{REGCV2} is an equivalent alternative; the EPRI grid-forming
\model{REGF2} (VSM) is the natural realism upgrade for a robustness study;
\model{REGCA1} (grid-following) and \model{REGF1}/\model{REGF3} (droop/dVOC) are
unsuitable because they do not expose a schedulable inertia constant.
\end{abstract}

\section{The requirement: what VIS actually schedules}
VIS treats the emulated inertia $M$ and damping $D$ of each grid-forming IBR as
\emph{decision variables}, alongside generation dispatch, under frequency-security
limits. The physical target being emulated is the synchronous-machine swing
equation (per unit, $\omega$ in p.u.):
\begin{equation}
2H\,\frac{d\omega}{dt} = P_m - P_e - D(\omega-1),
\label{eq:swing}
\end{equation}
whose immediate consequence, for an imbalance $\Delta P$, is the
rate-of-change-of-frequency
\begin{equation}
\dot f_0 = \frac{f_0\,\Delta P}{2H} = \frac{f_0\,\Delta P}{M},\qquad M \equiv 2H.
\label{eq:rocof}
\end{equation}
A suitable model must therefore (i) be \emph{grid-forming} (a controlled voltage
source that sets frequency, not a current source that follows it), and (ii)
expose an explicit \emph{inertia constant} $M{=}2H$ (and damping $D$) as a
tunable parameter. Droop gain alone is not inertia: it shapes the steady-state
power split, not the $d\omega/dt$ transient that \eqref{eq:rocof} governs.

\section{The ANDES model landscape}
ANDES~1.9.3 provides seven models in the \texttt{RenGen} group. \Cref{tab:models}
classifies them by control paradigm and, crucially, by whether an inertia
constant is an explicit parameter.

\begin{table}[h]
\centering\small
\caption{Renewable/converter models in ANDES 1.9.3 (\texttt{RenGen} group).}
\label{tab:models}
\begin{tabular}{@{}lllc@{}}
\toprule
Model & Control paradigm & Origin & Schedulable $M{=}2H$?\\
\midrule
\model{REGCA1} & Grid-\emph{following} (current src.) & WECC/PSS/E REGC\_A & \textcolor{red}{no}\\
\model{REGCP1} & Grid-forming (power/droop) & ANDES & \textcolor{red}{no ($M$)}\\
\model{REGCV1} & \textbf{Grid-forming VSG (swing eq.)} & ANDES (CURENT) & \textcolor{green!45!black}{\textbf{yes ($M$,$D$)}}\\
\model{REGCV2} & Grid-forming VSG (lag inner loop) & ANDES & \textcolor{green!45!black}{yes ($M$,$D$)}\\
\model{REGF1} & Grid-forming droop & EPRI GFM & \textcolor{red}{no}\\
\model{REGF2} & Grid-forming VSM & EPRI GFM & partial ($m_f$)\\
\model{REGF3} & Grid-forming dVOC & EPRI GFM & \textcolor{red}{no}\\
\bottomrule
\end{tabular}
\end{table}

\section{What REGCV1 actually solves}
\model{REGCV1} is documented as a ``voltage-controlled VSC with VSG control:
double-loop PI control and swing-equation-based VSG control.'' Its dynamic core
(from the ANDES source) is, for the virtual rotor speed deviation
$\Delta\omega$ and angle $\delta$,
\begin{align}
M\,\frac{d\,\Delta\omega}{dt} &= P_{\mathrm{ref}}' - P_e - D\,\Delta\omega,
& P_{\mathrm{ref}}' &= P_{\mathrm{ref}} - k_\omega\,\Delta\omega, \label{eq:regcv1_swing}\\
\frac{d\delta}{dt} &= 2\pi f_n\,\Delta\omega, & \omega &= 1 + \Delta\omega,
\end{align}
with an inner double-loop PI voltage/current controller
($V_d,V_q\!\to\!I_d,I_q$ through $r_a,x_s$) closing the electrical dynamics. The
ANDES parameters carry the exact physical meaning we need: \model{M} = ``emulated
startup time constant ($M{=}2H$)'' (\si{\second}) and \model{D} = ``emulated
damping coefficient'' (p.u.), with $k_\omega$ the active-power (speed) droop.

Equation~\eqref{eq:regcv1_swing} is \eqref{eq:swing} with $M{=}2H$: at the instant
of an imbalance, $d\Delta\omega/dt|_{0^+}=\Delta P/M$, reproducing
\eqref{eq:rocof}. Thus \emph{the scheduler's decision variable $M$ is literally
the inertia constant in the simulator's swing equation}, and the surrogate learns
the physically-grounded map $M\!\uparrow\,\Rightarrow\,|\dot f|\!\downarrow$. We
verified this causally in the pipeline: with a fixed disturbance, raising the
IBR inertia from $M{=}1$ to $M{=}7$ reduced $|\mathrm{RoCoF}|$ by $1.45\times$,
consistent with \eqref{eq:rocof}.

\section{Why not the alternatives}
\begin{description}[leftmargin=1.4em,style=nextline]
\item[\model{REGCA1} (grid-following).] A current source synchronised by a PLL; it
has no swing state and no inertia parameter (its parameters are current-limit and
low-voltage logic). It \emph{cannot} provide schedulable inertia and is therefore
unusable for VIS. It is, however, the right \emph{grid-following baseline} if one
later wants a GFM-vs-GFL comparison.
\item[\model{REGF1} (droop), \model{REGF3} (dVOC).] Grid-forming but inertia-free:
\model{REGF1} exposes droop gains ($w_{drp},Q_{drp}$), \model{REGF3} an oscillator.
Neither has an $M{=}2H$ parameter, so ``scheduling virtual inertia'' is undefined
for them. They address a different control question (droop/oscillator sharing),
not inertia allocation.
\item[\model{REGF2} (EPRI VSM).] The strongest alternative: a virtual synchronous
machine with an inertia-like coefficient $m_f$, damping $d_d$, a PLL, and explicit
current/power limits ($P_{\max},Q_{\max}$, \texttt{PQFLAG}). It is more
standards-aligned (EPRI/UNIFI GFM) and more realistic. We nonetheless retain
\model{REGCV1} as the primary model because: (a)~its $M{=}2H$ maps
\emph{directly and cleanly} to the VIS decision variable and to \eqref{eq:swing},
whereas $m_f$ requires an extra mapping; (b)~its smaller state set and absence of
hard current-limit switching keep the exact ReLU--MILP surrogate embedding
cleaner; and (c)~it preserves comparability with the foundational VIS work of
She~\emph{et al.} (2024), from the same CURENT/ANDES group, which uses this model.
\model{REGF2} is the appropriate model for a future \emph{realism robustness
study}---in particular to add explicit converter current-limit dynamics as a
predicted metric.
\item[\model{REGCV2} (VSG, lag inner loop).] Physically equivalent to
\model{REGCV1} with the inner current PI replaced by lag transfer functions; it
carries the same $M,D$. It is an acceptable substitute; \model{REGCV1} is the
canonical, slightly simpler choice and the one used upstream.
\end{description}

\section{Recommendation}
\begin{enumerate}[nosep]
\item \textbf{Primary model: \model{REGCV1}.} Decision-variable match ($M{=}2H$,
$D$), swing-equation transparency \eqref{eq:regcv1_swing}$\equiv$\eqref{eq:swing},
clean embeddability, and comparability with the reference VIS literature. This is
the model used throughout the current pipeline and case files
(\texttt{ieee39\_full\_ibrs.xlsx}), and the physics has been validated on the 50\,Hz
base.
\item \textbf{Robustness study (future): \model{REGF2}.} Re-run a subset with the
EPRI VSM to confirm the conclusions are not model-specific and to introduce
converter current-limit dynamics.
\item \textbf{Baseline (optional): \model{REGCA1}.} For an explicit
grid-forming-vs-grid-following inertia comparison.
\end{enumerate}
\textbf{Bottom line:} no model change is needed before the cluster campaign;
\model{REGCV1} is the correct and defensible choice for virtual inertia
scheduling.

\begin{thebibliography}{9}
\small
\bibitem{she2024} B.~She, F.~Li, H.~Cui, J.~Wang, Q.~Zhang, R.~Bo,
``Virtual Inertia Scheduling (VIS) for Real-Time Economic Dispatch of
IBR-Penetrated Power Systems,'' \emph{IEEE Trans. Sustain. Energy}, 15(2), 2024.
\bibitem{andes} H.~Cui, F.~Li, K.~Tomsovic, ``Hybrid Symbolic-Numeric Framework
for Power System Modeling and Analysis,'' \emph{IEEE Trans. Power Syst.}, 36(2),
2021. ANDES documentation: \url{https://docs.andes.app}.
\bibitem{milano} F.~Milano, F.~D\"orfler, G.~Hug, D.~J.~Hill, G.~Verbi\v{c},
``Foundations and Challenges of Low-Inertia Systems,'' \emph{PSCC}, 2018.
\bibitem{unifi} UNIFI Consortium, ``Specifications for Grid-Forming Inverter-Based
Resources, v2,'' NREL, 2024.
\bibitem{epri} EPRI, ``Grid-Forming Inverter Modeling (REGF1/2/3),'' EPRI
Memorandum (as implemented in ANDES).
\end{thebibliography}

\end{document}
